\documentclass[11pt,a4paper]{article}
\usepackage{amsmath,amssymb,amsthm}
\usepackage[margin=2.5cm]{geometry}
\usepackage{hyperref}

\newtheorem{definition}{Definition}[section]
\newtheorem{theorem}[definition]{Theorem}
\newtheorem{proposition}[definition]{Proposition}
\newtheorem{remark}[definition]{Remark}
\newtheorem{conjecture}[definition]{Conjecture}

\title{Hyperseed Formalization:\\
  Why the Good Guys Will Usually Win}
\author{ProtomegaTron (automated formalization)}
\date{2026-09-18}

\begin{document}
\maketitle

\begin{abstract}
We formalize Ben Goertzel's ``Why the Good Guys Will Usually Win''
(2023-07-09) within the Hyperseed ontology. The article argues that
open-mindedness and cooperation have structural advantages over
closed-mindedness and domination. Key mappings: openness as fiber richness,
cooperation as fiber composition, closed-mindedness as thin/rigid fiber,
trust as fiber coherence, and the multiverse argument as base-space
sampling favoring rich-fiber configurations.
\end{abstract}

\tableofcontents

\section{Fiber richness: the structural advantage of openness}

\begin{theorem}[Openness = rich fiber --- claims 1--5, 21]
Open-minded systems have richer fiber --- more perspectives, more inference
paths, more sections of the bundle. This gives a structural advantage:
\[
  |F_{\text{open}}| > |F_{\text{closed}}| \implies
  |\text{Solutions}(F_{\text{open}})| > |\text{Solutions}(F_{\text{closed}})|
\]
Rich fiber means more ways to understand and respond to any given base
situation. Open systems explore more of possibility space, find more
solutions, discover more resources, adapt more readily, and recombine
ideas across domains.
\end{theorem}

\section{Fiber composition: cooperation}

\begin{theorem}[Cooperation composes fibers --- claims 6--10, 22, 25]
Cooperation composes fibers from multiple agents, creating combined fiber
richer than any individual:
\[
  F_{\text{coop}} = F_1 \otimes F_2 \otimes \cdots \otimes F_n
  \supset F_i \quad \forall i
\]
The composed fiber has emergent sections no individual agent possesses.
Trust is fiber coherence --- compatible fiber structures that enable smooth
composition. Distrust prevents composition.
\end{theorem}

\section{Thin fiber: fragility of closed systems}

\begin{proposition}[Closed-mindedness = thin fiber --- claims 11--15, 23, 26]
Closed systems have thin, rigid fiber covering only a narrow band of base:
\[
  \text{Support}(F_{\text{closed}}) \subset B_{\text{narrow}}
  \implies \text{perturbation outside } B_{\text{narrow}} \Rightarrow \text{failure}
\]
Maintaining closed-mindedness requires fiber pruning --- removing sections
that lead to dissonant conclusions. This costs resources and weakens the
fiber, creating a self-reinforcing fragility cycle.
\end{proposition}

\section{Base-space measure: the multiverse argument}

\begin{conjecture}[Rich fiber dominates in measure --- claims 16--20, 24]
Across the full base space, rich-fiber configurations occupy more volume
and attract more trajectories than thin-fiber configurations:
\[
  \mu(\{b \in B : F_{\text{rich}}(b) \text{ survives}\})
  > \mu(\{b \in B : F_{\text{thin}}(b) \text{ survives}\})
\]
This is a statistical argument, not a local guarantee. The structural
advantages compound over long time spans: openness tends to win in
the long run even when it loses in the short run.
\end{conjecture}

\section{Cross-references}

\begin{remark}[Connections]
\begin{itemize}
\item \textbf{A BGI Manifesto} (2023-11-25): open development as safer
  than corporate capture --- same fiber richness argument.
\item \textbf{Evolving Deeply Ethical AGI} (2026-01-06): non-dual
  motivation as open fiber structure.
\item \textbf{Why Everyone Dies Gets AGI All Wrong} (2025-10-01):
  hybrid architecture as multi-fiber constraint --- same cooperation
  principle.
\end{itemize}
\end{remark}

\end{document}
